\ifx\allfiles\undefined
%生成目录超链接，使用UTF8字符集
\documentclass[12pt,UTF8]{ctexart}


\usepackage{mathtools}%花括号括起的公式
\usepackage{enumerate}%列表
\usepackage{graphicx}%表格
\usepackage[table]{xcolor}%彩色表格
\usepackage{booktabs}%调整表格行高


\renewcommand\tabcolsep{50pt}
\renewcommand\arraystretch{2}

%文本框
\usepackage{tcolorbox}
\tcbuselibrary{skins,vignette,raster,listings,listingsutf8,theorems, breakable,magazine,fitting}

%页面设置
\usepackage{geometry}%设置页面
\usepackage{setspace}%设置局部行距

\usepackage{ctex}
\usepackage{CJK}

%图像处理
\usepackage{graphicx}
\usepackage{float}%允许浮动
\usepackage{caption}%设置标题
\usepackage{graphicx, subfig}

%附录
\usepackage[toc,page]{appendix}

%导言区
\title{中小微企业的信贷决策}

%不要页眉
\pagestyle{plain}
\bibliographystyle{plain}

\geometry{a4paper}%将纸张设置为A4大小

%修改文档结构层次的名称
\usepackage{titlesec}%chapter1修改为一
\titleformat{\part}{\centering\Huge\bfseries}{\thepart}{1em}{}
\renewcommand\thepart{\chinese{part}}

%subsubsubsection的使用
%\usepackage{titlesec}
%\titleclass{\subsubsubsection}{straight}[\subsection]
%
%\newcounter{subsubsubsection}[subsubsection]
%\renewcommand\thesubsubsubsection{\thesubsubsection.\arabic{subsubsubsection}}
%\renewcommand\theparagraph{\thesubsubsubsection.\arabic{paragraph}}
%
%\titleformat{\subsubsubsection}
%{\normalfont\normalsize\bfseries}{\thesubsubsubsection}{1em}{}
%\titlespacing*{\subsubsubsection}
%{0pt}{3.25ex plus 1ex minus .2ex}{1.5ex plus .2ex}
%
%\makeatletter
%\renewcommand\paragraph{\@startsection{paragraph}{5}{\z@}%
%	{3.25ex \@plus1ex \@minus.2ex}%
%	{-1em}%
%	{\normalfont\normalsize\bfseries}}
%\renewcommand\subparagraph{\@startsection{subparagraph}{6}{\parindent}%
%	{3.25ex \@plus1ex \@minus .2ex}%
%	{-1em}%
%	{\normalfont\normalsize\bfseries}}
%\def\toclevel@subsubsubsection{4}
%\def\toclevel@paragraph{5}
%\def\toclevel@paragraph{6}
%\def\l@subsubsubsection{\@dottedtocline{4}{7em}{4em}}
%\def\l@paragraph{\@dottedtocline{5}{10em}{5em}}
%\def\l@subparagraph{\@dottedtocline{6}{14em}{6em}}
%\makeatother
%
%\setcounter{secnumdepth}{4}

\begin{document}
\else
\fi
%分文件开始


\subsection{问题四中旅游规划系统的开发}
\subsubsection{旅游规划系统的设计}
%先说一下一个用户界面应该有的特点，然后介绍一下我们的旅游规划系统的界面特点、结构特点等等

对一个旅游规划系统而言，其用户界面应具有以下几点主要特点：

首先是清楚性。用户界面设计的最终目的是为了使有不同需求的不同的人与系统通过沟通和功能来进行交互，因此一定的引导、解释是不可或缺的。

其次是简明性。用户界面应当为用户提供更加便捷、快速的渠道。

其三是高效率性。用户界面就像车辆，能准确的带用户去他想要到的地方。一个良好的界面能够让执行用户需求功能更快。

最后是吸引性。更具吸引力的界面不仅使得交互变得愉快，会让用户体验的满意度得以提升，还会让用户期待使用它。

~\\

本次我们所设计的张家界旅游规划系统正是遵从上述的四个主要特点来进行规划设计的。


\begin{figure}[H]
	\centering
	\includegraphics[width=1\linewidth]{figure/q4System}
	\caption{张家界旅行规划系统示意图}
	\label{fig:q4system}
\end{figure}


我们的用户界面由四个部分组合而成，分别为基本参数、必去地点、开支、结果面板。结构清晰，操作便捷，排列整齐有序。

在色彩方面，我们采用了以暖色调为主的风格设计，让使用的用户在视觉上不会感觉过于单一，从而影响使用的体验感与满意度。

在操作方面，我们采用键盘输入方式与按钮选择方式相结合，给用户提供更加快捷的方式，让他们可以在最短的时间内，找到满足要求的最合适的旅游线路以及对应的消费方案。在此系统中，还将根据近年来网上各景点的人气排名\cite{ref2}，推荐必去景点供用户选择，当然用户也可以手动输入想去的景点。用户可以根据系统引导，从左到右，从上到下依次进行操作，最终会在结果面板中得到答案。

在开支方面，我们设计的系统作出了进一步的细化供用户根据自身情况，进行相应的选择与输入，并将交通、食宿、门票各方面所需总额清楚的计算展现出来，让用户进一步了解各方面开支占比，从而极大的提高用户的使用效率，方便用户在接下来的旅游中进行详细的规划。


\subsubsection{旅游规划系统的使用}
%介绍一下这玩意怎么使用

首先，用户需在基本参数表中输入游玩的天数、人数、每日安排的游玩时间、总预算的基本信息，系统由此给出总游玩时间，从而进行下一步的安排.

接下来，用户可以在右侧表格必去地点中，采用手动输入或者按钮筛选输入的方式，选择在此次旅行安排中必去的地点，若有剩余的时间与预算，其余地点将由系统根据人气排名、最大限度利用时间与预算，进行合适的规划。

然后用户在开支面板中，首先挑选旅程所采用的交通方式，并依此自动计算所需费用。

其次，用户在食宿方面，可以自行输入每日食宿预算费用，同样系统也将自动计算总额。随后，根据用户选择的旅游景点，系统会依据每个景区的门票费用，自动计算出所需总额。

\begin{figure}[htbp]
	\centering
	\begin{minipage}[t]{0.49\textwidth}
		\centering
		\includegraphics[width=6cm]{figure/q4Input1}
	\end{minipage}
	\begin{minipage}[t]{0.49\textwidth}
		\centering
		\includegraphics[width=6cm]{figure/q4Input2}
	\end{minipage}
\end{figure}


\begin{figure}[H]
	\centering
	\includegraphics[width=0.7\linewidth]{figure/q4Input3}
	\label{fig:q4Input3}
\end{figure}


最后，用户会在结果面板中读取，此次张家界旅行的满意度总分、总开支、余额、总游玩时间与游玩景点数，并在右侧的旅游路线建议框中给出对应的具体游览景点顺序与路线。用户可由此做参考。

\begin{figure}[H]
	\centering
	\includegraphics[width=1\linewidth]{figure/q4Output}
	\caption{结果面板及旅游路线建议示意图}
	\label{fig:q4Output}
\end{figure}


\newpage


%分文件结束
\ifx\allfiles\undefined
\end{document}
\fi